% Version of August 04, 2020
%   Major cleanup of macros
%   Removal of obsolete macros (such as consti for MPI-1)
%   All ``internal'' macros renamed with II prefix (for *really* internal)
%   All user commands documented in the instr.tex file
%
% Version of September 15, 2010
%
% These are the TeX/LaTeX definitions that should be used by authors
% of the MPI standard.  Other definitions that are used to implement
% special effects (such as the line numbers or the appearance of
% chapter headings) are defined in mpi-sys-macs.tex
%
%
% To improve indexing for the document, \markindextrue will cause each index
% entry to create a mark at the point where the index is made, as well as adding
% it to the relevant index. These are not intended for use in the document;
% they are intended for use in the indexing macros that should be used in
% the document.
\usepackage{ifthen}
\usepackage{xstring}
\usepackage{atveryend}
\usepackage{etoolbox}

\newif\ifmarkindex
\markindexfalse
% These add items to the indices, and provide a way to review the
% index entries in the document. They all suppress adding index
% entries for items that occur in the table-of-contents, while
% including items that appear in section titles.
\def\IImpiindex#1{\ifcontents\else\ifmarkindex\hbox{\textbf{Index:#1}}\fi\index{#1}\fi}
\def\IImpimainindex#1{\ifcontents\else\ifmarkindex\hbox{\textbf{\underline{Index:#1}}}\fi\index{#1|uu}\fi}
% This version indexes the item in the main index UNLESS it is in the annex
% (e.g., the binding summaries)
\def\IImpimainindexA#1{\ifmpiInAnnex\IImpiindex{#1}\else\IImpimainindex{#1}\fi}

%%%%%%%%%%%%%%%%%%%%%%%%%%%%%%%%%%%%%%%%%%%%%%%%%%%%%%%%%%%%%%%%%%%%%%%%%%%%%
%% New Macros to format MPI, C, and Fortran items, using a consistent
%% naming style
%% <basename>[short][main][index|skip]{name}
%% \mpiconst is a special case - it takes a LaTeX optional argument that
%% allows the selection of a smaller font, which is needed sometimes in
%% tables.  The usage is \mpiconst[n]{name}, where ``n'' is an integer
%% representing the number of step sizes down to use.
%% Note: as of this writing, there are no uses of short with main, so
%% this combination is not defined. short and skip are also not used,
%% and that variation is not defined. These can be added if needed in
%% the document.
\def\IIfontsize#1{%
\ifnum#1=0\small\else
 \ifnum#1=1\footnotesize\else
  \ifnum#1=2\scriptsize\else
   \ifnum#1=3\tiny\else
    \ifnum#1>3\tiny\fi
   \fi
  \fi
 \fi
\fi}
%% For each of these commands, we define internal routines for
%% formatting the entry and for indexing the entry. All of the
%% variations of the macros make use of these

%% Functions and Parameters to functions

\newcommand{\IImpifuncfmt}[1]{\textsf{#1}}

\newcommand{\mpimaketmplabel}[1]{%
  \saveexpandmode\noexpandarg%
  \StrSubstitute{#1}{\_}{_}[\tmplabel]%
  \restoreexpandmode
}

% like \hyperlink but creates an MPI symbol label from the first argument
% format: \mpihyperlink{<symbol>}{<link text>}
\newcommand{\mpihyperlink}[2]{%
  \mpimaketmplabel{#1}%
  \hyperlink{\tmplabel}{#2}%
  \let\tmplabel\undefined%
}

% like \hypertarget but creates an MPI symbol label from the first argument
% format: \mpihypertarget{<symbol>}{<link text>}
\newcommand{\mpihypertarget}[2]{%
  \mpimaketmplabel{#1}%
  \hypertarget{\tmplabel}{#2}%
  \let\tmplabel\undefined%
}


% MPI API functions ONLY.  E.g., MPI_SEND.
% Names should be in all uppercase, which is the rule for the
% language-independent routine definitions
\newcommand{\mpifunc}[1]{\protect\gb\protect\mpihyperlink{#1}{\IImpifuncfmt{#1}}\IImpiindex{#1}}
\newcommand{\mpifuncnolink}[1]{\protect\gb\protect\IImpifuncfmt{#1}\IImpiindex{#1}}
\newcommand{\mpifuncskip}[1]{\protect\gb\protect\IImpifuncfmt{#1}}
\newcommand{\mpifuncmain}[1]{\mpihypertarget{#1}{\IImpifuncfmt{#1}}\IImpimainindex{#1}}
\newcommand{\mpifuncshort}[1]{\hbox{\mpihyperlink{#1}{\IImpifuncfmt{#1}}}\IImpiindex{#1}}
\newcommand{\mpifuncindex}[1]{\IImpiindex{#1}}
\newcommand{\mpifuncmainindex}[1]{\IImpimainindex{#1}}

% C or Fortran functions - not MPI API functions
\newcommand{\cfunc}[1]{\protect\gb\IImpifuncfmt{#1}}
\newcommand{\ffunc}[1]{\protect\gb\IImpifuncfmt{#1}}

% Parameters to MPI API functions
\newcommand{\mpiargshort}[1]{\IImpifuncfmt{#1}}
\newcommand{\mpiarg}[1]{\protect\gb\mpiargshort{#1}}

% Use these for parameters to C or Fortran functions respectively,
% esspecially MPI C and Fortran routine bindings.
\newcommand{\carg}[1]{\code{#1}}
\newcommand{\farg}[1]{\code{#1}}

%% Starter scripts like mpirun and mpiexec
\newcommand{\mpilaunch}[1]{\code{#1}\mpitermindex{#1}}
\newcommand{\mpilaunchskip}[1]{\code{#1}}
\newcommand{\mpilaunchmain}[1]{\code{#1}\IItermMainIndex{#1}}

%% Internal macros for MPI objects and definitions. Note that these
%% format the index entries so they have the same appearance as they
%% do in the text.
\newcommand{\IIconstfmt}[2][0]{\textsf{\IIfontsize{#1}#2}}
\newcommand{\IImpifmt}[1]{\textsf{#1}}
\newcommand{\IIconstidx}[2][CONST]{\IImpiindex{#1:#2@\IIconstfmt{#2}}}
\newcommand{\IIconstmainidx}[2][CONST]{\IImpimainindex{#1:#2@\IIconstfmt{#2}}}

% MPI Datatypes - e.g., MPI_DOUBLE
% Note this is a change from the previous macros - to use the same
% font size (\small) as the other constants.
\newcommand{\mpidtype}[1]{\protect\gb\mpidtypeshort{#1}}
\newcommand{\mpidtypeindex}[1]{\IIconstidx[DATATYPE]{#1}}
\newcommand{\mpidtypemain}[1]{\protect\gb\mpihypertarget{#1}{\IIconstfmt{#1}}\IIconstmainidx[DATATYPE]{#1}}
\newcommand{\mpidtypemainshort}[1]{\protect\mpihypertarget{#1}{\IIconstfmt{#1}}\IIconstmainidx[DATATYPE]{#1}}
\newcommand{\mpidtypeshort}[1]{\protect\mpihyperlink{#1}{\IIconstfmt{#1}}\IIconstidx[DATATYPE]{#1}}
\newcommand{\mpidtypeskip}[1]{\IIconstfmt{#1}}

% MPI C types - e.g., MPI_Comm, MPI_Aint or TYPE(MPI_Status)
\newcommand{\mpictype}[1]{\protect\gb\mpictypeshort{#1}}
\newcommand{\mpictypeshort}[1]{\protect\IIconstfmt{#1}\IIconstidx{#1}}
\newcommand{\mpictypemain}[1]{\protect\gb\protect\IIconstfmt{#1}\IIconstmainidx{#1}}
\newcommand{\mpictypeindex}[1]{\IIconstidx{#1}}
\newcommand{\mpictypemainindex}[1]{\IIconstmainidx{#1}}
\newcommand{\mpictypeskip}[1]{\protect\gb\IIconstfmt{#1}}
\newcommand{\mpiftype}[1]{\IIconstfmt{#1}\IIconstidx{#1}}

% MPI Fortran Integer kind reference.  Use \mpiftypekind{OFFSET} for
% INTEGER(KIND=MPI_OFFSET_KIND); this will index the use under
% MPI_OFFSET_KIND instead of under INTEGER(...)
% Note that this is new for MPI 4.0, and will need to be updated in
% the document.
\newcommand{\mpiftypekindshort}[1]{\code{INTEGER(KIND=MPI\_#1\_KIND)}%
\IIconstidx{MPI\_#1\_KIND}}
\newcommand{\mpiftypekind}[1]{\protect\gb\mpiftypekindshort{#1}}
\newcommand{\mpiftypekindmain}[1]{\code{INTEGER(KIND=MPI\_#1\_KIND)}%
\IIconstmainidx{MPI\_#1\_KIND}}

% MPI Error classes (which are also codes), including MPI_SUCCESS
% Also allows an optional size specifier
% Error codes are special in that some have main definitions a chapter
% while others don't. We provide a weak macro that sets a hyperref target
% if there is no main definition anywhere in the text.

% read in error label information from previous run
\IfFileExists{./\jobname.mel}{\input{\jobname.mel}}{}

% open file to write error label information to
\newwrite\mpierrorlabels
\immediate\openout\mpierrorlabels=\jobname.mel
% make sure the file is closed at the very end
\AtVeryEndDocument{%
    \closeout\mpierrorlabels%
}


\makeatletter
\newcommand{\error}[2][0]{\protect\mpimaketmplabel{#2}\protect\hyperlink{\tmplabel}{\errorskip[#1]{#2}}\IIconstidx{#2}}
%\ne`wcommand{\errormain}[2][0]{%
\DeclareRobustCommand{\errormain}[2][0]{%
    \mpimaketmplabel{#2}%
    \hypertarget{\tmplabel}{\errorskip[#1]{#2}}\IIconstmainidx{#2}%
    \csxdef{typeof@\tmplabel}{main}%
    \expandafter\ifcsname defined@\tmplabel\endcsname\relax\else%
        % mark as known
        \csxdef{defined@\tmplabel}{1}%
    \fi%
    \immediate\write\mpierrorlabels{\expandafter\csxdef{typeof@\tmplabel}{main}}%
}
\DeclareRobustCommand{\errorweak}[2][0]{%
	%\def\tmplabel{UNDEDINED}
	\begingroup
    \mpimaketmplabel{#2}%
    \expandafter\ifcsname defined@\tmplabel\endcsname%
        % we have seen this label aleady
        \error{#2}%
    \else%
		\csxdef{defined@\tmplabel}{1}%
        \expandafter\ifcsname typeof@\tmplabel \endcsname%
            \ifthenelse{\equal{\expandafter\csname typeof@\tmplabel\endcsname}{weak}}{%
                % weak definition: no other main definition has occured in a previous run and we have seen a weak definition
                \hypertarget{\tmplabel}{\errorskip[#1]{#2}}\IIconstmainidx{#2}%
            }{% some other main definition has occured in a previous run
                \error{#2}%
            }%
        \else % no definition yet
            \hypertarget{\tmplabel}{\errorskip[#1]{#2}}\IIconstmainidx{#2}%
            \csxdef{typeof@\tmplabel}{weak}%
        \fi%
		% write a marker that we have a weak label; may be overruled by a weak definition
        \immediate\write\mpierrorlabels{\expandafter\csxdef{typeof@\tmplabel}{weak}}%
	\fi%
	\endgroup
}
\makeatother
\newcommand{\errorskip}[2][0]{\protect\gb\protect\IIconstfmt[#1]{#2}}

% Predefined functions, such as MPI_COMM_DUP_FN.  These are not intended
% to be called by users, and are described in the standard as constants
% Note: For backward compatibility with the MPI 4.0 and earlier standards,
% these are added to the function index.  That should be revisited, and
% the MPI Function indexed reserved for functions that users should call,
% not these constants that are technically functions.
\newcommand{\mpiconstfunc}[1]{\protect\gb\protect\mpihyperlink{#1}{\IImpifuncfmt{#1}}\IIconstidx{#1}\mpifuncindex{#1}}
\newcommand{\mpiconstfuncskip}[1]{\protect\gb\IImpifuncfmt{#1}}
\newcommand{\mpiconstfuncmain}[1]{\protect\gb\mpihypertarget{#1}{\IImpifuncfmt{#1}}\IIconstmainidx{#1}\mpifuncmainindex{#1}}
\newcommand{\mpiconstfuncindex}[1]{\IIconstidx{#1}\mpifuncindex{#1}}


% Other constants.  Allows optional size reduction, as
% in \const[1]{name} (for 1 size smaller than the default)
\newcommand{\mpiconst}[2][0]{\protect\gb\protect\mpiconstshort[#1]{#2}}
\newcommand{\mpiconstmain}[2][0]{\protect\gb\mpiconstmainshort[#1]{#2}}
\newcommand{\mpiconstshort}[2][0]{\mpihyperlink{#2}{\IIconstfmt[#1]{#2}}\IIconstidx{#2}}
\newcommand{\mpiconstmainshort}[2][0]{\mpihypertarget{#2}{\IIconstfmt[#1]{#2}}\IIconstmainidx{#2}}
\newcommand{\mpiconstskip}[2][0]{\protect\gb\IIconstfmt[#1]{#2}}
% mpiconstmainindex adds *only* a main index entry. No text will appear in the
% text
\newcommand{\mpiconstmainindex}[1]{\mpihypertarget{#1}{\IIconstmainidx{#1}}}

%% C and Fortran types (e.g., double and INTEGER) are not indexed.
%% Do NOT use this for types such as MPI_Aint - use mpictype instead
%% (and mpiftype for the Fortran counterparts, such as TYPE(MPI_Status))
\newcommand{\ctype}[1]{\protect\gb\texttt{#1}}
\newcommand{\ftype}[1]{\protect\gb\texttt{#1}}
\newcommand{\ctypeshort}[1]{\texttt{#1}}
%% This special version of short ensures that the Fortran names are
%% not hyphenated. Consider making this the default for single-word
%% Fortran types.
\newcommand{\ftypeshort}[1]{\texttt{\hbox{#1}}}

%% Callbacks (called typedef in the macros used for MPI-1 through MPI-3).
% These are for the definitions of Callback routines, such as the
% attribute copy and delete callbacks
% Language-independent names
% They are also used to index the parameters to the functions that
% create them; use mpicallback for the language-independent and
% mpiccallback for the C-language callbacks (and ask for a Fortran
% version mpifcallback if needed)
% (need to disambiguate that in the future)
% TODO: use \IImpifmt here?
\newcommand{\IIcallbackfmt}[1]{\protect\gb\protect\textsf{#1}}
\newcommand{\mpicallback}[1]{\protect\gb\protect\mpihyperlink{#1}{\textsf{#1}}\mpicallbackindex{#1}}
\newcommand{\mpicallbackmain}[1]{\protect\gb\protect\mpihypertarget{#1}{\textsf{#1}}\mpicallbackmainindex{#1}}
\newcommand{\mpicallbackskip}[1]{\protect\gb\textsf{#1}}
\newcommand{\mpicallbackindex}[1]{\IImpiindex{TYPEDEF:#1@\textsf{#1}}}
\newcommand{\mpicallbackmainindex}[1]{\IImpimainindex{TYPEDEF:#1@\textsf{#1}}}

%% C definition (typedef) for callbacks
% As a function name, uses the default size, not \small, for the text.
% These are the MAIN (defining) index entries for these callbacks
\newcommand{\mpiccallbackdefindex}[1]{\IImpimainindexA{TYPEDEF:#1@\textsf{#1}}}
%% Fortran definition for callbacks
\newcommand{\mpifcallbackdefindex}[1]{\IImpimainindexA{TYPEDEF:#1@\textsf{#1}}}

%% References to language-specific callbacks
\newcommand{\mpiccallback}[1]{\code{#1}\mpicallbackindex{#1}}
\newcommand{\mpifcallback}[1]{\code{#1}\mpicallbackindex{#1}}
\newcommand{\mpifcallbackskip}[1]{\code{#1}}


%% String constants in various contexts
% Internal routines to format {foo} as "foo", and to index under foo
\newcommand{\IIstringidx}[2][CONST]{\IImpiindex{#1:#2@\protect\IIfmatstring{#2}}}
\newcommand{\IIstringmainidx}[2][CONST]{\IImpimainindex{#1:#2@\protect\IIfmatstring{#2}}}

% Info is for MPI_Info predefined strings.  \infokey{keyname} and
% \infoval{keyvaluename}
% We use a slightly different const index to indicate that these are
% info strings.
% Note that these format as strings, but must be used without the quotes.
% E.g., \infokey{cb\_size}, not \infokey{"cb\_size"}
\newcommand{\infokey}[1]{\infokeyskip{#1}\IIstringidx[INFOCONST=KEY]{#1}}
\newcommand{\infoval}[1]{\infovalskip{#1}\IIstringidx[INFOCONST=VAL]{#1}}
\newcommand{\infovalshort}[1]{\IIfmatstring{#1}\IIstringidx[INFOCONST=VAL]{#1}}
\newcommand{\infokeymain}[1]{\infokeyskip{#1}\IIstringmainidx[INFOCONST=KEY]{#1}}
\newcommand{\infovalmain}[1]{\infovalskip{#1}\IIstringmainidx[INFOCONST=VAL]{#1}}
\newcommand{\infokeyskip}[1]{\protect\gb\IIfmatstring{#1}}
\newcommand{\infovalskip}[1]{\protect\gb\IIfmatstring{#1}}

% Miscellaneous strings, including (when fixed in the document)
% external32 and native for data representations.
\newcommand{\mpistring}[1]{\protect\gb\mpistringshort{#1}}
\newcommand{\mpistringshort}[1]{\IIfmatstring{#1}\IIstringidx[STRINGCONST]{#1}}
\newcommand{\mpistringskip}[1]{\protect\gb\IIfmatstring{#1}}
\newcommand{\mpistringmain}[1]{\IIfmatstring{#1}\IIstringmainidx[STRINGCONST]{#1}}

%% -- Backward compatability.  The uses of these should be fixed in
%% the document, and then these can be removed.

\newcommand{\mpiskipfunc}[1]{\protect\gb\IImpifuncfmt{#1}}
% Note that this macro is not the same as mpifuncshort (and the use is
% dubious in the document and should be re-examined)
\newcommand{\mpishortfunc}[1]{\hbox{\IImpifuncfmt{#1}}}
\newcommand{\mpishortarg}[1]{\mpiargshort{#1}}

% const is now mpiconst to emphasize that it should ONLY be used with
% MPI constants
\newcommand{\const}[1]{\mpiconst{#1}}
\newcommand{\constmain}[1]{\mpiconstmain{#1}}
\newcommand{\constshort}[1]{\mpiconstshort{#1}}
\newcommand{\constshortmain}[1]{\IIconstfmt[0]{#1}\IIconstmainidx{#1}}
\newcommand{\constskip}[1]{\mpiconstskip{#1}}

% Use the \const[2]{name} form instead for the different sized const.
% Note that all versions of mpiconst support the size option.
\newcommand{\constsmall}[1]{\mpiconst[1]{#1}}
\newcommand{\constfootnotesize}[1]{\mpiconst[2]{#1}}
\newcommand{\constscriptsize}[1]{\mpiconst[3]{#1}}
\newcommand{\consttiny}[1]{\mpiconst[3]{#1}}

% Use mpidtype versions instead
\newcommand{\type}[1]{\mpidtype{#1}}
\newcommand{\typeindex}[1]{\mpidtypeindex{#1}}
\newcommand{\typemain}[1]{\mpidtypemain{#1}}

% Use infokeyskip instead (or infovalskip, if an info value)
\newcommand{\infoskip}[1]{\infokeyskip{#1}}

% Note that cdeclindex is sometimes misused in the document - ensure
% that it is replaced with the correct macro
\newcommand{\cdeclindex}[1]{\mpictypeindex{#1}}
\newcommand{\cdeclmainindex}[1]{\mpictypemainindex{#1}}

% Add an index for ``BigCount'' C routines
\newcommand{\MPIbindindex}[1]{\IImpimainindex{BCFUNCTION:#1}}
\newcommand{\MPIbindindexuse}[1]{\IImpiindex{BCFUNCTION:#1}}
% Use this for the ``BigCount'' C routines in text.  Using \mpifunc is
% incorrect, and \mpifunc is only for the language-neutral versions
\newcommand{\MPIBCfunc}[1]{\protect\gb\protect\IImpifuncfmt{#1}\IImpiindex{BCFUNCTION:#1}}

% Make descriptions use the same indent as paragraphs
\setlist[description]{leftmargin=\docparindent}

%%%%%%%%%%%%%%%%%%%%%%%%%%%%%%%%%%%%%%%%%%%%%%%%%%%%%%%%%%%%%%%%%%%%%%%%%%%%%

% fancyvrb defines Verbatim, which is a slightly better Verbatim environment
% Regrettably, the key feature needed, commandchars, does not work correctly
% in our environment (it doesn't accept arbitrary grouping characters).
% See README-2.2 for instructions on using Verbatim along with the above
% update macros.
\usepackage{fancyvrb}
%
% Place some penalty for doing the break
% The penalty for a ``\gb'' should be greater than a \hyphenpenalty.
% \hyphenpenalty is 50 in plain.tex.
% Using more than 8em for the flexible glue can cause bad breaks where
% that much space isn't needed.  See vlgb as an option
\def\gb{\penalty10000\hskip 0pt plus 8em\penalty4800\hskip 0pt plus-8em%
\penalty10000}
\def\vlgb{\penalty10000\hskip 0pt plus 16em\penalty4800\hskip 0pt plus-16em%
\penalty10000}
%%%%%%%%%%%%%%%%%%%%%%%%%%%%%%%%%%%%%%%%%%%%%%%%%%%%%%%%%%%%%%%%%%%%%%%%%%%%%

%
% Define an example environment, which provides numbers that can be accessed
% by using \label{eq-name} ... \ref{eq-name}.
\makeatletter
\DeclareNewTOC[listname={List of Examples},type=mpiexample]{loex}
\DeclareTOCStyleEntry[numwidth=2.5em]{tocline}{mpiexample}
\newcommand{\mpiexamples}[1]{%
\addxcontentsline{loex}{mpiexample}[\thechapter.\theexample]{#1}}
\newcounter{example}[chapter]

%
% Define a counter that is used to ensure page references are correct
% when applied to unnumbered environments including advice and user.
% Examples are already correct becaue they have an associated counter
\newcounter{pagerefs}
% Use pagelabel instead of label when you intend to refer to this label
% for a page number AND the label is not intended for a (sub)*section.
% Note that the advice/user/rationale environments set things up so
% that a regular \label works.
% The \stepcounter ensures that the page number is correct
% The \phantomsection is needed so that hyperlinks go to the correct page,
% rather than the most recent section
\newcommand{\pagelabel}[1]{\stepcounter{pagerefs}\phantomsection\label{#1}}

\usepackage[skins,breakable]{tcolorbox}

\newtcolorbox{examplebox}{%
colback=gray!0, %
opacityfill=0, %
opacityframe=1, %
enhanced jigsaw, %
boxrule=1.5pt, %
boxsep=0pt, %
left=5pt, %
right=4pt, %
top=5pt, %
bottom=5pt, %
sharp corners=northwest, %
grow to left by=5pt, %
grow to right by=10pt, %
breakable=true, %
before={\ifvmode\else\par\fi\vspace{5pt}}, %
after skip=5pt, %
colframe=lightgray %
}

\newif\ifInExample
\InExamplefalse
\newenvironment{example}{\addtocounter{example}{1}\InExampletrue\noindent\begin{examplebox}\textbf{Example
    \arabic{chapter}.\arabic{example}.} \interlinepenalty6000%
\edef\@currentlabel{\arabic{chapter}.\arabic{example}}%
% Phantomsection is needed so that hyperref will know to jump here if this
% example has a \label
\phantomsection
\label{example-\thechapter.\theexample}%
%\rm
}{\end{examplebox}\InExamplefalse}
\makeatother
%
% Use \exindex{MPI\_FUNC} to generate an index entry for MPI
% functions/constants
%\newcommand{\exindex}[1]{\relax}
% FIXME: exindex should only be used within an example, and examples should all
% be in the example environment
%\newcommand{\exindex}[1]{\ifInExample\IImpiindex{EXAMPLES:#1}\else\errmessage{exindex used outside of an example}\fi}
\newcommand{\exindex}[1]{\IImpiindex{EXAMPLES:#1}\ifInExample\else\typeout{Warning: exindex used outside of an example}\fi}
%
% New indexing for indices
% Use \Exindex[shorttitle]{language}{index entry}{namelist}
% as in
% \Exindex{C}{Use of send and receive}{MPI\_Send,MPI\_Recv}
% The index entry is used as the short title if no short title is given.
% The name list may contain any mixture of:
%   MPI Function names
%   MPI Declarations (e.g., MPI_Aint), expressed as d:name, e.g., d:MPI_Aint
%   MPI Constants (e.g., MPI_COMM_WORLD), expressed as c:name
%   MPI Callback Typedefs (e.g., MPI_Handler_function), expressed as t:name
%   MPI Infokeys or values, expressed as i:name, e.g., i:same_op
\newif\ifexindexerror
\exindexerrorfalse
\newcommand{\Exindex}[4][]{\ifInExample\IImpiindex{NEWEXAMPLES:\arabic{chapter}.\arabic{example}::#2::#3::#4}\def\mytst{#1}\mpiexamples{\ifx\mytst\empty{#3}\else{#1}\fi}\else\ifexindexerror\errmessage{Exindex used outside of an example}\else\IImpiindex{NEWEXAMPLES:NA::#2::#3::#4}\typeout{Warning: Exindex used outside of an example}\fi\fi}

%%%%%%%%%%%%%%%%%%%%%%%%%%%%%%%%%%%%%%%%%%%%%%%%%%%%%%%%%%%%%%%%%%%%%%%%%%%%%
%
% The following commands make it eaiser to create a list of participants
% in a table, which is used in the frontmatter.  Use like this:
% \begin{tabular}{lll}
% first name \EE
% second name \EE
%...
% name 1007 \EE
% last name in list \FlushEntries
% \end{tabular}
% This avoids having to keep moving & and \\ markers around as a list is
% constructed and updated.
% This is setup for 3 columns.  Redefine \colmax for the number of columns
% if different than 3 is desired.
% For especially long names, use both multicolumn and \EEi:
% ,,, \multicolumn{2}{l}{A long name}\EEi\EE ...
% The \EEi advances the counter that tracks the number of columns used.
% If this is a common occurance, we can define a macro that combines these,
% such as \EElongname{name} -> \multicolumn{2}{l}{#1}\EEi
\newcounter{IIcolidx}
\setcounter{IIcolidx}{0}
\def\IIcolmax{3}
% \EE for ``EndEntry''
\def\EE{\stepcounter{IIcolidx}\ifnum\value{IIcolidx}=\IIcolmax\setcounter{IIcolidx}{0}\\\else&\fi}
\def\EEi{\stepcounter{IIcolidx}}
% Use \FlushEntries *instead* of \EE after the last name.
%\def\FlushEntries{\stepcounter{IIcolidx}\ifnum\value{IIcolidx}=\IIcolmax\\\else&\FlushEntries\fi\setcounter{IIcolidx}{0}}
% This simplified version issues an end-of-row if a row has started, and
% resets the counter IIcolidx.
% It does *not* generate the missing alignment chars (&) but this does
% not seem to be a problem in MacTeX 2025.
\def\FlushEntries{\ifnum\value{IIcolidx}=0\else\\\fi\setcounter{IIcolidx}{0}}

%% %
%% % Here's a version that will work for equal spaced columns, and could handle
%% % an occasional long line
%% \def\EEbx{\vbox\bgroup\hbox to 2in\bgroup}
%% \def\EEb{\noindent\EEbx}
%% \def\EE{\hfil\egroup\egroup\ \EEbx}
%% \def\FlushEntries{\hfil\egroup\egroup}

%
%
%%%%%%%%%%%%%%%%%%%%%%%%%%%%%%%%%%%%%%%%%%%%%%%%%%%%%%%%%%%%%%%%%%%%%%%%%%%%%
% Ensure that the formatting of MPI names is consistent.  These use
% \MPIname/ (with the trailing slash) to avoid the common problem of a missing
% space, as in ``\MPI foo'' which goes to ``MPIfoo'' if \MPI is defined
% as ``MPI''

\def\MPI/{\textsf{MPI}}              % MPI macro
\def\mpi/{\textsf{MPI}}              % MPI macro
\def\MPII/{\textsf{MPI-1}}           % MPI-1 macro -- can't easily do MPI1
\def\mpii/{\textsf{MPI-1}}           % MPI-1 macro
\def\MPIIDOTO/{\textsf{MPI-1.0}}     % MPI-1.0 macro - NOTE it has a O (Oh)
                                     % not 0 (zero)
\def\mpiidoto/{\textsf{MPI-1.0}}     % MPI-1.0 macro - NOTE it has a O (Oh)
                                     % not 0 (zero)
\def\MPIIDOTI/{\textsf{MPI-1.1}}     % MPI-1.1 macro
\def\mpiidoti/{\textsf{MPI-1.1}}     % MPI-1.1 macro
\def\MPIIDOTII/{\textsf{MPI-1.2}}    % MPI-1.2 macro
\def\mpiidotii/{\textsf{MPI-1.2}}    % MPI-1.2 macro
\def\MPIIDOTIII/{\textsf{MPI-1.3}}   % MPI-1.3 macro
\def\mpiidotiii/{\textsf{MPI-1.3}}   % MPI-1.3 macro
\def\MPIII/{\textsf{MPI-2}}          % MPI-2 macro -- can't easily do MPI2
\def\mpiii/{\textsf{MPI-2}}          % MPI-2 macro
\def\MPIIIDOTO/{\textsf{MPI-2.0}}    % MPI-2.0 macro - NOTE it has a O (Oh)
                                     % not 0 (zero)
\def\mpiiidoto/{\textsf{MPI-2.0}}    % MPI-2.0 macro - NOTE it has a O (Oh)
                                     % not 0 (zero)
\def\MPIIIDOTI/{\textsf{MPI-2.1}}    % MPI-2.1 macro
\def\mpiiidoti/{\textsf{MPI-2.1}}    % MPI-2.1 macro
\def\MPIIIDOTII/{\textsf{MPI-2.2}}   % MPI-2.2 macro
\def\mpiiidotii/{\textsf{MPI-2.2}}   % MPI-2.2 macro
\def\MPIIII/{\textsf{MPI-3}}         % MPI-3 macro -- can't easily do MPI3
\def\mpiiii/{\textsf{MPI-3}}         % MPI-3 macro
\def\MPIIIIDOTO/{\textsf{MPI-3.0}}   % MPI-3.0 macro
\def\mpiiiidoto/{\textsf{MPI-3.0}}   % MPI-3.0 macro
\def\MPIIIIDOTI/{\textsf{MPI-3.1}}   % MPI-3.1 macro
\def\mpiiiidoti/{\textsf{MPI-3.1}}   % MPI-3.1 macro
\def\MPIIV/{\textsf{MPI-4}}          % MPI-4 macro
\def\mpiiv/{\textsf{MPI-4}}          % MPI-4 macro
\def\MPIIVDOTO/{\textsf{MPI-4.0}}    % MPI-4.0 macro
\def\mpiivdoto/{\textsf{MPI-4.0}}    % MPI-4.0 macro
\def\MPIIVDOTI/{\textsf{MPI-4.1}}    % MPI-4.1 macro
\def\mpiivdoti/{\textsf{MPI-4.1}}    % MPI-4.1 macro
\def\MPIV/{\textsf{MPI-5}}           % MPI-5 macro
\def\mpiv/{\textsf{MPI-5}}           % MPI-5 macro
\def\MPIVDOTO/{\textsf{MPI-5.0}}     % MPI-5.0 macro
\def\mpivdoto/{\textsf{MPI-5.0}}     % MPI-5.0 macro
\def\MPIVDOTI/{\textsf{MPI-5.1}}     % MPI-5.1 macro
\def\mpivdoti/{\textsf{MPI-5.1}}     % MPI-5.1 macro
\def\MPIJOD/{\textsf{MPI-JOD}}       % MPI-JOD macro
\def\mpijod/{\textsf{MPI-JOD}}       % MPI-JOD macro
\def\MPIRT/{\textsf{MPI/RT}}         % RT MPI macros.
%
\def\RMA/{\textsf{RMA}}              % RMA macro
%%%%%%%%%%%%%%%%%%%%%%%%%%%%%%%%%%%%%%%%%%%%%%%%%%%%%%%%%%%%%%%%%%%%%%%%%%%%%
% In some cases, we need to indicate future versions of MPI.  These macros
% allow the chapter authors to indicate that version symbolically.
% Before producing a releasable version of the MPI document, these macros
% *must* be replaced with the specific MPI versions intended.
\def\MPInext/{\textsf{MPI-next}}
\def\MPInextoh/{\textsf{MPI-next-major}}
%%%%%%%%%%%%%%%%%%%%%%%%%%%%%%%%%%%%%%%%%%%%%%%%%%%%%%%%%%%%%%%%%%%%%%%%%%%%%

% Vertical space used to offset the function definitions, constant
% lists, and mpicodeblocks
\newlength{\IIcodeSpace}
\setlength{\IIcodeSpace}{.4cm}

%% Common leader for function definitions
\def\IIfuncdefleader{
    \vspace{\IIcodeSpace}
    \vspace{\IIcodeSpace}
    \noindent
%    \samepage % Samepage is obsolete.  And do we want it?
     \encouragesamepage  % Encourage function defintion to be on
                         % one page, but allow it to cross a page if
                         % the definition is long.  This uses TeX
                         % interline penalties (samepage essentially
                         % sets those penalties to infinity)
    \raggedright
    \parskip0pt
    \hangindent 7em\hangafter=1
}
\newenvironment{funcdef}[2]{
    \IIfuncdefleader
    \mpimaketmplabel{#1}%
	\hypertarget{\tmplabel}{\IImpifuncfmt{#1#2}}\mpifuncmainindex{#1#2}
    \MPIfunclist
}{\end{list} \vspace{\IIcodeSpace}\gdef\IIbindingcurlang{\IIbindingdef}}

%
% This is a hack to force a linebreak
\newenvironment{funcdef2}[2]{
    \IIfuncdefleader
    \mpimaketmplabel{#1}%
	{\hypertarget{\tmplabel}{\IImpifuncfmt{#1}}\flushline
    \IImpifuncfmt{#2}}\mpifuncmainindex{#1#2}
    \MPIfunclist
}{\end{list} \vspace{\IIcodeSpace}\gdef\IIbindingcurlang{\IIbindingdef}}

% Special funcdef for functions with no arguments
\newenvironment{funcdefna}[2]{
    \vspace{\IIcodeSpace}
    \vspace{\IIcodeSpace}
    \noindent
    \mpimaketmplabel{#1}%
	\hypertarget{\tmplabel}{\IImpifuncfmt{#1#2}}\mpifuncmainindex{#1#2}\par
}{\vspace{\IIcodeSpace}\gdef\IIbindingcurlang{\IIbindingdef}}


% Based on the definition of \samepage in LaTeX, but with smaller penalties
\makeatletter
\def\encouragesamepage{\interlinepenalty4000
\postdisplaypenalty3000
\interdisplaylinepenalty3000
\@itempenalty3000
\@beginparpenalty3000
\@endparpenalty3000
}
\makeatother
                                        % see page 77, the TeX book.
\newcommand{\funcarg}[3]{\item[\hbox to 45pt{\textsf{#1} \hfill} \mpiarg{#2}\hfill]{\parbox[t]{3.2in}{\raggedright\small #3}}}

%%%%%%%%%%%%%%%%%%%%%%%%%%%%%%%%%%%%%%%%%%%%%%%%%%%%%%%%%%%%%%%%%%%%%%%%%%%%%
% makeindex considers a double-quote (") a special character, and escaping
% it with the default character (backslash) confuses either makeindex or
% LaTeX, depending on exactly how used. We use indirection with protect
% to keep \dq from being expanded until after the line is processed by
% makeindex.
% We also use the \tt font for the double quote to ensure an ``unsigned''
% double quote - completely vertical ticks.
\def\dq{{\small\ttfamily\char`"}}
\newcommand{\IIfmatstring}[1]{\IIconstfmt{\protect\dq{}#1\protect\dq}}
%%%%%%%%%%%%%%%%%%%%%%%%%%%%%%%%%%%%%%%%%%%%%%%%%%%%%%%%%%%%%%%%%%%%%%%%%%%%%
\newcommand{\IN}[0]{\IIconstfmt{IN}}
\newcommand{\OUT}[0]{\IIconstfmt{OUT}}
\newcommand{\INOUT}[0]{\IIconstfmt{INOUT}}

%%%%%%%%%%%%%%%%%%%%%%%%%%%%%%%%%%%%%%%%%%%%%%%%%%%%%%%%%%%%%%%%%%%%%%%%%%%%%
%
% This defines a list environment for constants, along with macros for the
% constants in the list.
\newenvironment{constlist}[1][160pt]{
    % lists already have their vertical space here
    \noindent
    \begin{list}{}{                     % see pg 113 of Lamport's book
        \setlength{\leftmargin}{#1}
        \addtolength{\leftmargin}{40pt}  % 40pt is common indent in other parts of standard
        \setlength{\labelwidth}{#1}
        \setlength{\labelsep}{10pt}
        \setlength{\itemindent}{10pt}
        \setlength{\itemsep}{-5pt}
        \setlength{\topsep}{-5pt}
        \setlength{\listparindent}{0pt}
    }
}{\end{list} \vspace{\IIcodeSpace}}

\newcommand{\noconstitem}[2]{\item[#1\hfill]{\noindent #2}}
\newcommand{\constitem}[2]{\item[\const{#1}\hfill]{\noindent #2}}
\newcommand{\constmainitem}[2]{\item[\constmain{#1}\hfill]{\noindent #2}}
\newcommand{\constitemtwo}[3]{\item[\const{#1}, \const{#2}\hfill]{\noindent #3}}
\newcommand{\constitemthree}[4]{\item[\const{#1}, \const{#2}, \const{#3}\hfill]{\noindent #4}}
\newcommand{\constitemonly}[1]{\constitem{#1}{\ }}

% mpidtypeitem is like constitem, but for MPI datatypes
\newcommand{\mpidtypeitem}[2]{\item[\mpidtype{#1}\hfill]{\noindent #2}}
\newcommand{\mpidtypeitemonly}[1]{\mpidtypeitem{#1}{\ }}
\newcommand{\mpidtypemainitem}[2]{\item[\mpidtypemain{#1}\hfill]{\noindent #2}}
%%%%%%%%%%%%%%%%%%%%%%%%%%%%%%%%%%%%%%%%%%%%%%%%%%%%%%%%%%%%%%%%%%%%%%%%%%%%%
%
% Code Formatting
%  code    - Used for code in one of the language bindings (e.g., C)
%  mpicode - Used for language-independent code
\def\code#1{\texttt{#1}}
\def\mpicode#1{\textsf{#1}}

% Formatting for legend to the table in A.2 Summary of the Semantics of all
% operations-related MPI procedures.
% Note that the text uses this formatting but the table itself does not use
% any formatting
\def\mpilegend#1{\textsf{#1}}

%%%%%%%%%%%%%%%%%%%%%%%%%%%%%%%%%%%%%%%%%%%%%%%%%%%%%%%%%%%%%%%%%%%%%%%%%%%%%
%
% Use mpiterm when introducing a term that you want emphasized and indexed.
% Use mpitermni for terms that should not be indexed (ni = not indexed)
% Place the index first to allow emph to (possibly) correct spacing.
% mpitermdefindex produces only an index entry. The combination of ni and
% index allows index entries that differ from the text.
% References to section titles are marked with titleindex.
%
% An attempt was made to detect whether a term was the first use.  However,
% the presense of \_ in some of the terms caused the simple code (using
% \csname #1 \endcsname) to fail, and code to sanitize the argument is
% complex.  So we fell back to the simple choice here.
% QUERY: This sort of check is more easily handled during index
% processing, and could be built into buildindices
\def\IItermIndex#1{\IImpiindex{TERM:#1}}
\def\IItermMainIndex#1{\IImpimainindex{TERM:#1}}

\def\mpiterm#1{\IItermIndex{#1}\emph{#1}}
\def\mpitermni#1{\emph{#1}}
\def\mpitermindex#1{\IItermIndex{#1}}
\def\mpitermdef#1{\textbf{#1}\IItermMainIndex{#1}}
\def\mpitermdefni#1{\textbf{#1}}
\def\mpitermdefindex#1{\IItermMainIndex{#1}}
% Sometimes a term needs a more complex index entry; e.g., a subentry
\def\mpitermsub#1#2{\IItermIndex{#2}\emph{#1}}

% Special macro for lb\_marker and ub\_marker that are a term
% which is a very special category and should be printed with sf
\def\mpiublb#1{\textsf{#1}\IItermIndex{#1}}

\def\mpitermtitleindex#1{\IImpiindex{TERM:#1|bold}}
\def\mpitermtitleindexsubmain#1#2{\IImpiindex{TERM:#1 #2|bold}\IImpiindex{TERM:#2!#1|bold}}
% e.g. \mpitermtitleindexsubmain{Point-to-Point}{Communication}
% results in:     Point-to-Point Communication, 23
%                 Communication,
%                    Point-to-Point, 23
\def\mpitermtitleindexmainsub#1#2{\IImpiindex{TERM:#1!#2|bold}\IImpiindex{TERM:#2|bold}}
% e.g. \mpitermtitleindexmainsub{Message}{Envelope}
% results in:     Message
%                    Envelope, 27
%                 Envelope, 27

%%%%%%%%%%%%%%%%%%%%%%%%%%%%%%%%%%%%%%%%%%%%%%%%%%%%%%%%%%%%%%%%%%%%%%%%%%%%%
%
% Use flushline to force a linebreak without right justifying the line.
\def\flushline{\hfill\hbox{}\linebreak}
%
%%%%%%%%%%%%%%%%%%%%%%%%%%%%%%%%%%%%%%%%%%%%%%%%%%%%%%%%%%%%%%%%%%%%%%%%%%%%%
% Drawn from HPF Forum
%
% Use the pagerefs counter and a phantomsection to ensure that (a) the
% page number in a \pageref is correct and (b) use phantomsection to
% ensure that a hyperlink goes to this spot instead of the beginning of the
% (sub)*section

\newenvironment{rationale}{\begin{list}{}{\setlength{\leftmargin}{\docparindent}}\item[]\textit{Rationale.}%
\stepcounter{pagerefs}\phantomsection}{{\textrm{(\textit{End of rationale.})}} \end{list}}

\newenvironment{implementors}{\begin{list}{}{\setlength{\leftmargin}{\docparindent}}\item[]\textit{Advice
        to implementors.}\stepcounter{pagerefs}\phantomsection%
}{{\textrm{(\textit{End of advice to implementors.})}} \end{list}}

\newenvironment{users}{\begin{list}{}{\setlength{\leftmargin}{\docparindent}}\item[]\textit{Advice to users.}%
\stepcounter{pagerefs}\phantomsection}{{\textrm{(\textit{End of advice to users.})}} \end{list}}

%%%%%%%%%%%%%%%%%%%%%%%%%%%%%%%%%%%%%%%%%%%%%%%%%%%%%%%%%%%%%%%%%%%%%%%%%%%%%%
%macros for language binding: mpibind, mpifnewbind, mpifbind, and fargs:
% The binding headings can be turned on and off
% with \bindingheadingtrue and \bindingbheadingfalse. This is useful
% when all bindings are for a single language, and the heading is not needed.
\newif\ifbindingheading
\bindingheadingtrue
% The \ifvmode ... \fi adds a sort of good vertical break that allows a
% page to break right abouve the bindingheading, and will provide some
% blank space at the bottom of the page.
\newcommand{\IIbindingheading}[1]{\ifbindingheading{%
\ifvmode\vskip 0pt plus 100pt\penalty4800\vskip 0pt plus-100pt\fi
\noindent\textbf{#1}\par\nobreak}\fi}

\def\IIbindingcname{C binding}
\def\IIbindingfname{Fortran binding}
\def\IIbindingfnamespecial{Fortran binding (the following procedure is not available with mpif.h)}
\def\IIbindingfnewname{Fortran 2008 binding}

% This if controls whether there is a heading
\newif\ifbindinglabel
\bindinglabelfalse

% Because of how these are used in the \if tests, these definitions
% need to be a single character.
% Special case: ``s'' is really a subset of f
\def\IIbindingc{c}
\def\IIbindingf{f}
\def\IIbindingfspecial{s}
\def\IIbindingfnew{F}
\def\IIbindingdef{N}
\def\IIbindingcurlang{\IIbindingdef}
\def\IIbindingissue#1#2{\ifbindinglabel\if\IIbindingcurlang#1\else
\if\IIbindingfspecial#1\gdef\IIbindingcurlang{\IIbindingf}\else
\gdef\IIbindingcurlang{#1}\fi#2\fi\fi}
%% Use the following to reset the binding language to ``NEW'' binding
%% - E.g., A new binding will generate a heading in the correct
%% language, even if the last binding was in the same language.  Use,
%% for example, when all bindings are in C - but blocks of bindings
%% are separated in the text.
\newcommand{\bindinglangreset}{\gdef\IIbindingcurlang{N}}
%% Use \bindingislisttrue to turn *off* the spacing after the binding body
%% This is needed when there are consecutive binding blocks with no
%% intervening text. This occurs in the bindings chapter for the list
%% of handle conversion routines.
\newif\ifbindingislist \bindingislistfalse
% The widow and club penalty attempt to avoid page breaks after the first or
% before the last line of the binding. Note that these are local (within a
% set of braces) and may be overridden by many other settings, including
% the behavior of line breaks. That's a long way of saying that these
% might help but do not guarantee that pages don't break at bad spots in
% the binding body.
% The parskip0pt is needed for bindings within an itemize (e.g., in the
% deprecated interfaces chapter). Otherwise, there are extra spaces added
% between some lines (ones that look like paragraph breaks).
% Note that in an itemize environment, the hangindent doesn't work.
\newcommand{\IIbindingbody}[1]{%
{\raggedright \hangindent 7em\hangafter=1\ttfamily
\parskip0pt
\widowpenalty10000\clubpenalty10000
#1 \par \ifbindingislist\else\medskip\fi}}

%%%%%%%%%%%%%%%%%%%%%%%%%%%%%%%%%%%%%%%%%%%%%%%%%%%%%%%%%%%%%%%%%%%%%%%%%%%%%%
% IMPORTANT!
% Any changes to the binding macros must be coordinated with changes
% to the binding tools (in binding-tool), buildapplang, and buildindices
%%%%%%%%%%%%%%%%%%%%%%%%%%%%%%%%%%%%%%%%%%%%%%%%%%%%%%%%%%%%%%%%%%%%%%%%%%%%%%

\newcommand{\mpibind}[2]{{\IIbindingissue\IIbindingc{\IIbindingheading{\IIbindingcname}}%
\IIbindingbody{int \mpihyperlink{#1}{#1}#2}}}
\newcommand{\mpibindmain}[2]{{\IIbindingissue\IIbindingc{\IIbindingheading{\IIbindingcname}}%
\IIbindingbody{int \mpihypertarget{#1}{#1}#2}}}

\newcommand{\mpibindnotint}[3]{{\IIbindingissue\IIbindingc{\IIbindingheading{\IIbindingcname}}%
\IIbindingbody{#1 \mpihyperlink{#2}{#2}#3}}}
\newcommand{\mpibindnotintmain}[3]{{\IIbindingissue\IIbindingc{\IIbindingheading{\IIbindingcname}}%
\IIbindingbody{#1 \mpihypertarget{#2}{#2}#3}}}


%special macro for not including C binding in index
%% UNUSED
%\newcommand{\mpiskipbind}[1]{{\IIbindingissue\IIbindingc{\IIbindingheading{\IIbindingcname}}%
%\IIbindingbody{int #1}}}

% special macro that avoids the int in front
% should be used for C functions only that want to be in index
%% UNUSED
%\newcommand{\mpiemptybind}[2]{{\IIbindingissue\IIbindingc{\IIbindingheading{\IIbindingcname}}%
%\IIbindingbody{#2 #1}}}
%
% This version indexes the name
% \mpiemptybindidx{routine}{returnvalue}{argument}{indexed name}
\newcommand{\mpiemptybindidx}[4]{{\ifvmode\else\par\fi\medskip
\IIbindingissue\IIbindingc{\IIbindingheading{\IIbindingcname}}%
\IIbindingbody{\mpihypertarget{#4}{#3 #1#2\mpifuncmainindex{#4}}}}}
%
% This version does not index the name, but has same argument list
%% Used only in appLang-CNames.tex, generated ONLY by MAKE-APPLANG
\newcommand{\mpiemptybindNOidx}[3]{{\IIbindingissue\IIbindingc{\IIbindingheading{\IIbindingcname}}%
\IIbindingbody{#3 #1#2}}}

% special macro that avoids the int in front
% should be used for routines where don't want in index
%% UNUSED
%\newcommand{\mpiskipemptybind}[2]{{\IIbindingissue\IIbindingc{\IIbindingheading{\IIbindingcname}}%
%\IIbindingbody{#2 #1}}}

% special macro for typedef int in front of SUBROUTINE in C
%
% Note that this is indexed.

\newcommand{\mpitypedefbindmain}[2]{{\IIbindingbody{typedef int \mpihypertarget{#1}{#1}#2;\mpiccallbackdefindex{#1#2}}}}
\newcommand{\mpitypedefbindvoidmain}[2]{{\IIbindingbody{typedef void \mpihypertarget{#1}{#1}#2;\mpiccallbackdefindex{#1#2}}}}
\newcommand{\mpitypedefbind}[2]{{\IIbindingbody{typedef int \mpihyperlink{#1}{#1}#2;\mpicallbackindex{#1#2}}}}
\newcommand{\mpitypedefbindvoid}[2]{{\IIbindingbody{typedef void \mpihyperlink{#1}{#1}#2;\mpicallbackindex{#1#2}}}}
%
% special macro for callback prototype definitions in Fortran
\newcommand{\mpifsubbindmain}[2]{{\IIbindingbody{SUBROUTINE \mpihypertarget{#1}{#1}#2}\mpicallbackmainindex{#1}}}
\newcommand{\mpifsubbind}[2]{{\IIbindingbody{SUBROUTINE \mpihyperlink{#1}{#1}#2}\mpicallbackindex{#1}}}
%
% Use \tt because we format more than a single paragraph, and \texttt
% will complain
\newcommand{\mpifnewsubbind}[2]{{\ttfamily\noindent ABSTRACT INTERFACE\par\noindent\quad\IIbindingbody{SUBROUTINE #1#2}\mpifcallbackdefindex{#1}}}

% special macro for normal SUBROUTINE and FUNCTION definitions in Fortran
% Fortran function names are the same as the MPI LIS names, so we reference those
\newcommand{\mpifbind}[2]{{\IIbindingissue\IIbindingf{\IIbindingheading{\IIbindingfname}}%
\IIbindingbody{\mpihyperlink{#1}{#1}#2}}}
% the main macro does not define a new hyperref label
\newcommand{\mpifbindmain}[2]{{\IIbindingissue\IIbindingf{\IIbindingheading{\IIbindingfname}}%
\IIbindingbody{#1#2}}}

% Fortran functions returning a value
\newcommand{\mpifbindret}[3]{{\IIbindingissue\IIbindingf{\IIbindingheading{\IIbindingfname}}%
\IIbindingbody{#1 \mpihyperlink{#2}{#2}#3}}}
% the main macro does not define a new hyperref label
\newcommand{\mpifbindretmain}[3]{{\IIbindingissue\IIbindingf{\IIbindingheading{\IIbindingfname}}%
\IIbindingbody{#1 #2#3}}}


% special macro for MPI_Status_f082f | f2f08, which uses the
% "Fortran binding (the following procedure is not available with mpif.h)" label
\newcommand{\mpifbindspecial}[2]{{\IIbindingissue\IIbindingfspecial{\IIbindingheading{\IIbindingfnamespecial}}%
\IIbindingbody{#1#2}}}

% Fortran 2008 bindings have the same name as C bindings so we reference the C bindings
% and the main macro does not define a new target
\newcommand{\mpifnewbind}[2]{{\IIbindingissue{\IIbindingfnew}{\IIbindingheading{\IIbindingfnewname}}%
\IIbindingbody{\mpihyperlink{#1}{#1}#2}}}
\newcommand{\mpifnewbindmain}[2]{{\IIbindingissue{\IIbindingfnew}{\IIbindingheading{\IIbindingfnewname}}%
\IIbindingbody{#1#2}}}
% of course there are exceptions: bindings without C bindings must declare a target...
\newcommand{\mpifnewbindmainref}[2]{{\IIbindingissue{\IIbindingfnew}{\IIbindingheading{\IIbindingfnewname}}%
\IIbindingbody{\mpihypertarget{#1}{#1}#2}}}
% special case for Fortran functions returning a value (e.g., MPI_WTIME)
\newcommand{\mpifnewbindret}[3]{{\IIbindingissue{\IIbindingfnew}{\IIbindingheading{\IIbindingfnewname}}%
\IIbindingbody{#1 \mpihyperlink{#2}{#2}#3}}}
\newcommand{\mpifnewbindretmain}[3]{{\IIbindingissue{\IIbindingfnew}{\IIbindingheading{\IIbindingfnewname}}%
\IIbindingbody{#1 #2#3}}}

% special macro for text instead of a binding; used for routines that were deprecated until MPI-2.2:
\newcommand{\mpifnewnonebind}[1]{{\IIbindingbody{\textrm{#1}}}}

% special macro to skip fortran binding in index
%% UNUSED
%\newcommand{\mpifskipbind}[1]{{\IIbindingbody{#1}}}

% The \\* creates an unbreakable linebreak, which helps avoid page breaks
% just after the Fortran function arg list (since fargs is only used for
% the Fortran bindings)
\def\fargs{\\*\advance\leftskip 2em\hangindent5em\hangafter=1}
%
\def\fnarg{\par\hangindent5em\hangafter=1}
\def\fnfinalarg{\par\nobreak\hangindent5em\hangafter=1}

%% Use \bindindent/ to end a line and start a new line with a consistent
%% indent.  This replaces the recommendation to use
%%  \\\ \ \ \   (that last \ is followed by a space)
%% when continuing a declaration line.  This is mostly used in the Fortran
%% bindings.
%% The 7em indent is the same as for the c bindings - 5+2 (2 already for
%% fargs)
\def\bindindent/{\\\hskip5em}

% Mark the annex so that macros can adjust behavior in the annex.
% This is used, for example, to select which index to use for items
% that appear in the document and in the Annex
% Originally defined for overloaded Fortran mpi module routines
% that should only be printed in the Annex A
\newif\ifmpiInAnnex \mpiInAnnexfalse
\newcommand{\mpifoverloadStartInAnnex}{\mpiInAnnextrue}
% Replace \\[3pt] with \par\vspace{3pt}
\newcommand{\mpifoverloadOnlyInAnnex}[1]{{\ifmpiInAnnex
\ifvmode\else\par\fi\vspace{3pt} {\raggedright \hangindent
7em\hangafter=1\textrm{If the Fortran compiler provides \ftype{TYPE(C\_PTR)}, then the above is overloaded by:}
\vspace{-7pt}\ttfamily \begin{tabbing} mm\=mm\=mm\=mmmm\= \kill #1 \end{tabbing} \vspace{-7pt}}\fi}}

%% - The following version is needed when removing the \MPIupdate{3.0}... macros:
%% \newcommand{\mpifoverloadOnlyInAnnexA}[1]{{\ifmpiInAnnex
%% \\[3pt] {\raggedright \hangindent
%% 7em\hangafter=1\rm If the Fortran compiler provides \ftype{TYPE(C\_PTR)}, then overloaded by:
%% \vspace{-7pt}\tt \begin{tabbing} mm\=mm\=mm\=mmmm\= \kill #1 \end{tabbing} \vspace{-7pt}}\fi}}

%%%%%%%%%%%%%%%%%%%%%%%%%%%%%%%%%%%%%%%%%%%%%%%%%%%%%%%%%%%%%%%%%%%%%%%%%%%%%
% Tell TeX about special words in the standard so that they can be
% hyphenated automatically
\hyphenation{in-ter-com-mun-i-cat-or in-tra-com-mun-i-cat-or}
\hyphenation{in-ter-com-mun-i-ca-tion in-tra-com-mun-i-ca-tion}
\hyphenation{sem-a-phore}
\hyphenation{drop-ped}
\hyphenation{to-po-lo-gies}
\hyphenation{con-tig-u-ous buf-fers}
\hyphenation{car-te-sian}

% Use foo\hyphen/barfiddle instead of foo-barfiddle to allow TeX to hyphenate
% barfiddle (by default, TeX won't hyphenate any words that contain a hyphen)
\def\hyphen/{\hskip0pt\relax-\penalty0\hskip0pt\relax}
%%%%%%%%%%%%%%%%%%%%%%%%%%%%%%%%%%%%%%%%%%%%%%%%%%%%%%%%%%%%%%%%%%%%%%%%%%%%%


%%%%%%%%%%%%%%%%%%%%%%%%%%%%%%%%%%%%%%%%%%%%%%%%%%%%%%%%%%%%%%%%%%%%%%%%%%%%%
%% Adjust figure spacing - the default is 0pt + 1fil at the top, which will
%% put figures that take up a part of a page but not the full page sort of
%% in the middle.  This fixes that.
\makeatletter
\setlength{\@fptop}{0pt}
\makeatother
%%%%%%%%%%%%%%%%%%%%%%%%%%%%%%%%%%%%%%%%%%%%%%%%%%%%%%%%%%%%%%%%%%%%%%%%%%%%%
%%
%% Referring to a section or other named object in the document
%% Removed pagereference.  Consider making this a configuration option
%\newcommand{\namedref}[2]{#1~\ref{#2} on page~\pageref{#2}}
\newcommand{\namedref}[2]{#1~\ref{#2}}
\newcommand{\sectionref}[1]{\namedref{Section}{#1}}
\newcommand{\annexref}[1]{\namedref{Annex}{#1}}
% Define section range to include the page ranges if you really insist on
% cluttering up the text with information useful only to the printed version
\newcommand{\sectionrange}[2]{Sections~\ref{#1} to~\ref{#2}}
%%%%%%%%%%%%%%%%%%%%%%%%%%%%%%%%%%%%%%%%%%%%%%%%%%%%%%%%%%%%%%%%%%%%%%%%%%%%%
%%
%% Language-independent code block environment
%% This needs to be used with care.  If it is used inline, follow it with
%% \noindent to keep TeX from starting a new paragraph.
\newenvironment{mpicodeblock}{\ifvmode\else\par\fi\vspace{\IIcodeSpace}%
\noindent\sffamily\quad}{\ifvmode\else\par\fi\vspace{\IIcodeSpace}}

%%
%% Use \XXX/ for a ``function name'' wildcard
\def\XXX/{\textsf{XXX}}

%%
%% In the RARE cases where it is necessary to refer explictly to a
%% chapter or section title, use this command if the section or title
%% does not have a valid label
\newcommand{\textTitleFont}[1]{\textrm{#1}}
%%
%% Where the section does have a label, use \nameref{label-value}.

%%%%%%%%%%%%%%%%%%%%%%%%%%%%%%%%%%%%%%%%%%%%%%%%%%%%%%%%%%%%%%%%%%%%%%%%%%%%%%

%%%%%%%%%%%%%%%%%%%%%%%%%%%%%%%%%%%%%%%%%%%%%%%%%%%%%%%%%%%%%%%%%%%%%%%%%%%%%
% In numerous places in the standard, an unnumbered heading/label is used
% to begin a paragraph.  Much (though not all) of the document has used
% \paragraph*{heading} for this purpose. This is problematic, because (a)
% paragraph is a sectioning command, below subsubsection and (b) the
% formatting of the heading is weak, so that these headings do not stand
% out. Rather than use lower-level font commands, we have defined a
% special command for these paragraph headings, \paraheading. Some other
% places should use the appropriate level of sectioning command (e.g.,
% typically \subsubsection).
%%%%%%%%%%%%%%%%%%%%%%%%%%%%%%%%%%%%%%%%%%%%%%%%%%%%%%%%%%%%%%%%%%%%%%%%%%%%%
\newcommand{\paraheading}[1]{\vspace{3.0ex}\noindent\textbf{#1}}

% Text Blocks
% In some cases, there are common blocks of text needed in the standard
% The following macros define these blocks of text
%%
%% Used for the MPI_T interface to standardize the language in the
%% language-independent description.
\newcommand{\textoutargs}[2]
{
\funcarg{\OUT}{#1}{buffer to return the string containing #2 (string)}
\funcarg{\INOUT}{{#1}\_len}{length of the string and/or buffer for \mpiarg{#1} (integer)}
}

\newcommand{\textoutdesc}[2]
{
The arguments \mpiarg{#1} and \mpiarg{{#1}\_len} are used to return #2
as described in Section~\ref{sec:mpit:strings}.%
}

\newcommand{\textoutoptinfo}[2]
{
\MPI/ can optionally return an info object containing the default hints set
for #2.  If the
argument to \mpiarg{#1} provided by the user is the \code{NULL} pointer,
this argument is ignored, otherwise an \MPI/
implementation is required to return all hints that are supported by the
implementation for #2 and have default values specified; any user-supplied
hints that were not ignored by the implementation; and any additional hints
that were set by the implementation. If no such hints exist, a handle to a
newly created info object is returned that contains no key/value pair. The
user is responsible for freeing \mpiarg{#1} via \mpifunc{MPI\_INFO\_FREE}.
}

% Use these for the world and sessions model names.
% WARNING: If followed by a space, write as \wpm{} or \spm{};
% otherwise the space will be lost (this is why the MPI macros
% are \MPI/ etc.)
\newcommand{\wpm}{World Model}
\newcommand{\spm}{Sessions Model}

\def\MPIisinitialized/{\MPI/ is \mpifuncindex{MPI\_INIT}\mpifuncindex{MPI\_INIT\_THREAD}initialized}
\def\MPIisfinalized/{\MPI/ is \mpifuncindex{MPI\_FINALIZE}finalized}

%%%%%%%%%%%%%%%%%%%%%%%%%%%%%%%%%%%%%%%%%%%%%%%%%%%%%%%%%%%%%%%%%%%%%%%%%%%%%
% This is a temporary fix to make add some space between deprecated items
% for use in the chapter on deprecated features.
\def\deprecatedsep{\ifvmode\else\par\fi\bigskip}

%%%%%%%%%%%%%%%%%%%%%%%%%%%%%%%%%%%%%%%%%%%%%%%%%%%%%%%%%%%%%%%%%%%%%%%%%%%%%
%% Macro to create a thick horizontal line in table environments
%% Use \thickhline instead of \hline if you need a thicker horizontal line
\makeatletter
\def\thickhline{%
  \noalign{\ifnum0=`}\fi\hrule \@height \thickarrayrulewidth \futurelet
   \reserved@a\@xthickhline}
\def\@xthickhline{\ifx\reserved@a\thickhline
               \vskip\doublerulesep
             \fi
      \ifnum0=`{\fi}}
\makeatother

\newlength{\thickarrayrulewidth}
\setlength{\thickarrayrulewidth}{2\arrayrulewidth}
%%%%%%%%%%%%%%%%%%%%%%%%%%%%%%%%%%%%%%%%%%%%%%%%%%%%%%%%%%%%%%%%%%%%%%%%%%%%%

%
% This defines a rule that can be used to underline an entire line.
% This is intended to be used in examples that have column headings (replacing
% the use of multiple hyphens, as in ----------)
\def\wideunderline{\hbox to 0pt{\raisebox{-2pt}[0pt][0pt]{\vbox to 0pt{\hrule height 0.4pt width \linewidth}}\hss}}

%%%%%%%%%%%%%%%%%%%%%%%%%%%%%%%%%%%%%%%%%%%%%%%%%%%%%%%%%%%%%%%%%%%%%%%%%%%%%%
% Nicer listings

\usepackage{listings}
%\usepackage{courier}      % basewidth=0.5em??
%\usepackage{lmodern}      % basewidth=0.6em
\usepackage{inconsolata}  % basewidth (not specified)

\definecolor{MPIlightgrey}{RGB}{240,240,240}
% Use this with basicstyle=\lstsmall when a slightly smaller font is required
% for the listing. For example,
% \begin{lstlisting}[language=C,basicstyle=\lstsmall]
\def\lstsmall{\footnotesize\ttfamily}

% Note: When listing keywords or other lists, if a keyword is a subsequence
% of another keyword, list the shortest one FIRST.  I.e.,
% ...,MPI_CHAR,MPI_CHARACTER,...

% Adjust the default style
\lstset{language=C,
        keepspaces=true,
        showstringspaces=false,
        basicstyle=\small\ttfamily,
        columns=fixed,
%        columns=fixed,basewidth=0.5em,
%        columns=fixed,basewidth=0.6em,
        backgroundcolor=\color{MPIlightgrey},
}

% Use kpsewhich lstlang1.sty to find to location of the definitions of the
% langauges known to lstlisting.  You need te consult this file to update
% certain choices, such as deleteing keywords (they can only be deleted
% from the level in which they are defined - so you need to check lstlang1.sty
% to find that level).

% Define an "MPI" dialect of default C.  This is required to get
% the morekeywords option to work, and is more in keeping with the design
% of the listing package.
\lstdefinelanguage[MPI]{C}[]{C}%
{%
        keywordsprefix={MPI_},%
% We must list any PMPI names
        morekeywords={PMPI_Init,PMPI_T_init_thread,PMPI_T_pvar_get_num,%
PMPI_T_pvar_get_info,PMPI_T_pvar_session_create,PMPI_T_pvar_handle_alloc,%
PMPI_T_pvar_start,PMPI_T_pvar_read,PMPI_T_pvar_handle_free,%
PMPI_T_pvar_session_free,PMPI_T_finalize,PMPI_Example,%
PMPI_Finalize,PMPI_Recv,PMPI_Send,%
% mpi info keys used in examples
mpi_hw_resource_type,
},%
% make #pragma a "keyword" so that it doesn't need to start in column 1,
% which is a requirement in the lstlisting package for directives
        otherkeywords={\#pragma},
}

\lstdefinelanguage[MPI]{Fortran}[90]{Fortran}%
{
        keywordsprefix={MPI_},%
        morekeywords={mpi_send},%
        % Yes, LEN and SIZE are defined in groups 2 and 3, and must
        % be deleted from both
        deletekeywords=[2]{LEN,SIZE,STAT,STATUS},
        deletekeywords=[3]{COUNT,INDEX,LEN,RANK,SHAPE,SIZE,SUM},
        mathescape=false
}

%% Add c_ptr - c_f_pointer is defined in the base 2008 dialect
\lstdefinelanguage[MPI08]{Fortran}[08]{Fortran}%
{
        keywordsprefix={MPI_},%
        morekeywords={PMPI_ISEND,c_ptr},%
        % Yes, LEN and SIZE are defined in groups 2 and 3, and must
        % be deleted from both
        deletekeywords=[2]{LEN,SIZE,STAT,STATUS},
        deletekeywords=[3]{COUNT,INDEX,LEN,RANK,SHAPE,SIZE,SUM},
        mathescape=false
}

% Use this environment around blocks of LaTeX that must be processed as a unit
% by tohtml (when creating an HTML version of the standard). It is a no-op
% for the PDF versoin.
\newenvironment{tohtml}{}{}

%% The following are special versions for Fortran examples that use
%% either a keyword as a variable, or a language extension. It does not
%% seem possible to apply these changes to a particular lstlisting
%% environment without defining a new language dialect
%% These are to be avoided if possible - for clarity, examples should
%% avoid using Fortran keywords, even though that is allowed in the
%% language

% A special version of MPI08 that includes CPP directives, a common
% extension in some Fortran compilers
% Also, the example that uses this dialect uses "function" as a user-defined
% function name. We remove it as a keyword
\lstdefinelanguage[MPI08dirs]{Fortran}[MPI08]{Fortran}%
{
        moredelim=*[directive]\#,%
        moredirectives={define,elif,else,endif,error,if,ifdef,ifndef,line,%
      include,pragma,undef,warning},%
        deletekeywords={FUNCTION}
}[keywords,comments,strings,directives]

% A special version of MPI that excludes IN and OUT, which are used in
% one example
\lstdefinelanguage[MPInoinout]{Fortran}[MPI]{Fortran}%
{
      deletekeywords={IN,OUT}%
}

% A special version of MPI that includes LEN, which is used as a Fortran
% keyword in one example
\lstdefinelanguage[MPI08withlen]{Fortran}[MPI08]{Fortran}%
{
      morekeywords={LEN}%
}

% HPF is used in a few places. This is very limited, and just covers what is
% needed in the example used in the MPI Standard
\lstdefinelanguage[HPF]{Fortran}[95]{Fortran}%
{
     morekeywords={PROCESSORS,DISTRIBUTE,CYCLIC,ONTO,BLOCK},%
     otherkeywords={!HPF$}%
}[keywords,comments,strings,directives]

%%%%%%%%%%%%%%%%%%%%%%%%%%%%%%%%%%%%%%%%%%%%%%%%%%%%%%%%%%%%%%%%%%%%%%%%%%%%%%
